\documentclass{ltxdoc}
% \documentclass[
% ngerman,
% paper=a4,
% asymmetric,
% fontsize=9pt,
% % headings=small,
% draft=false,
% parskip=half,
% % listof=totoc,         % verschiebt Minitoc/Hyperref
% % bibliography=totoc,   % verschiebt Minitoc/Hyperref
% % index=totoc,
% % noonelinecaption
% footnotes=multiple,
% toc=flat,
% ]{scrbook}
% \setlength{\parindent}{1.5em}
% \usepackage{scrhack}

% UTF-8 Schriftkodierung
\usepackage[utf8]{inputenc}\usepackage[T1]{fontenc}
% get rid of the "no room for new \ldots"
% standard TeX registers are limitied to 256, etex increases the max.
\usepackage{etex}
\usepackage{doc}
\usepackage{scrbase}
\newcounter{chapter}
\usepackage{aass-common}
\usepackage{aass-features}
\usepackage{aass-fonts}
\usepackage{aass-text}
\OnlyDescription
\begin{document}

\OnlyDescription
\section{Über das Project}

\section{\LaTeX}
\DocInput{aass-common.sty}

LaTeX2e, and most contributed macro packages, are now written in a literate programming style, with source and documentation in the same file. This format, known as ‘doc’, in fact originated before the days of the LaTeX project as one of the “Mainz” series of packages. A documented source file conventionally has the suffix .dtx, and will normally be ‘stripped’ before use with LaTeX; an installation file (.ins) is normally provided, to automate this process of removing comments for speed of loading. To read the comments, you can run LaTeX on the .dtx file to produce a nicely formatted version of the documented code. Several packages can be included in one .dtx file (they’re sorted out by the .ins file), with conditional sections, and there are facilities for indexes of macros, etc.

Anyone can


\end{document}